\paragraph{Credit accrual needs shared amounts before it can enter an exact book.}
The insurer examples use rational amounts. Dated credit accrual introduces a
different problem: separate numerical bounds can enclose each entry correctly
while losing the dependence that makes the journal balance. In the retained
July credit example, let $F=(11/10)^{181/365}$ be the dimensionless accrual
factor for the 181 days from 1 January to 1 July 2025. The signed entries are
$100(F-1)$ to gross carrying amount, zero to allowance, and $-100(F-1)$ to
interest income, all in the native currency units. They sum exactly to zero.
Selecting every entry's upper bound from the numerical report instead leaves
a residual of $10^{-16}$ currency units. A midpoint lies inside the factor's
enclosure but fails its exact defining equation. Numerical accuracy alone
therefore does not supply exact journal amounts.

One shared positive day factor preserves the dependence. For the supplied
fixed annual rate $r>-1$, set $a=1+r$ and choose the unique positive $z$ with
$z^{365}=a$. A $d$-day accrual uses $z^d$; discounting uses $z^{-d}$. For any
integer exponent $e$, write $e=365q+s$ with $0\leq s<365$. Then
\[
 z^e=(z^{365})^qz^s=a^qz^s,
 \qquad
 \sum_e c_ez^e=\sum_{s=0}^{364}b_sz^s .
\]
Here the rational coefficients $c_e$ are collected into rational coefficients
$b_s$ after the substitution. Negative $q$ is valid because $a>0$. Addition
and multiplication followed by this reduction preserve evaluation at the
same positive root. For monetary amounts the coefficients carry currency
units; the powers of $z$ are dimensionless. At zero rate, $z=1$ and the
expression reduces directly to a rational amount.

If every reduced coefficient is zero, equality is proved. A nonzero remainder
does not prove inequality, because $z^{365}-a$ need not be irreducible. For
example, the algebra-only choice $a=32$ gives $(z^{73})^5=32=2^5$, so positivity
implies $z^{73}=2$ even though $z^{73}-2$ has degree below 365. This example is
outside the retained financial rate range; it tests the limit of the equality
procedure. The reference returns an uncertified result whenever a required
identity does not reduce to zero. It provides division by rational constants
and inverse day powers, with no general polynomial inverse or ordering test.

The read-only \texttt{CreditLinkedSchedule/1} recomputes and binds each complete
retained request and native result before reconstructing the expressions.
Across nine schedules it preserves 57 stages and 63 journals, including
partial and full write-off, prepayment, negative and zero rates, leap dates
and irregular receipts. At each stage, net carrying amount equals gross less
allowance, gross amounts reconcile to their dated origins, remaining expected
cash discounts to net carrying amount, and cash plus net assets equals opening
capital plus profit. Accumulated signed entries reproduce the same balances
and income. Eighty independently authored targets and 652 exact identities
pass; 986 certified numerical enclosures lie within the original bounds. The
enclosures check consistency; the shared expressions establish the identities.

Complete exchanges and refusal evidence are retained under
{\small\path{data/accounting-master/}} in the directory
{\small\path{credit-linked-measurements-20261025/}}. Three
retained cases also project exactly to rational amounts. The existing rational
company book still requires a separate admission and account-mapping step.
The other six cases need an explicitly specified exact amount extension or
an authored quantization policy with its treatment of rounding differences.
The checked expressions establish neither choice and make no ledger mutation.
